
The transfer improves over the coarse portion of the tested sequences, but its composite round-trip error saturates for real fields on the finest pair. For the controlled analytic field, refining the \ROMS mesh while holding \MPAS fixed shows that the \ROMS characteristic length is not what limits the result, while cancellation between the forward and reverse legs makes the round trip a different quantity from the one-way forcing error. Which stage sets the floor remains unattributed, since we did not exercise the horizontal and vertical operators separately. Requiring the vertical interpolant to conserve layer integrals makes the composed tensor-product transfer first order overall, which is consistent with the measured decrease, though confirming the rate by fitting would require a longer sequence than the four grids used here.

The regional applications show finer-scale relative-vorticity structure and domain-dependent kinetic energy changes across open-ocean, boundary current, coastal, and estuarine settings. The four configurations sit at $\hRatio\approx0.21$ (North Atlantic), $0.24$ (Gulf Stream), $0.35$ (Bay of Bengal), and $0.06$ (Chesapeake Bay), so three of them lie near the flattening point identified in \S\ref{conv:sec:conv} while Chesapeake Bay is far below it, refined for reasons of estuarine geometry rather than transfer accuracy. As a result, the Chesapeake Bay case utilizes blended bathymetry information to produce better fidelity in regions where \MPAS can only provide sparse data.

Several limitations qualify some of these presented results. The \MPAS characteristic length is fixed and varies in space, while the \ROMS evaluation support changes as that mesh is refined (depending on the interior set, $\mathcal{I}$). Additionally, velocities are transferred between \MPAS and \ROMS models as componentwise scalars, and remain cell centered, without rotation to the \ROMS curvilinear basis, interpolation to staggered points, or any constraint on the discrete divergence. Future studies should explore use of Mimetic schemes to achieve edge-based, high-order consistent projections that do not require 3 levels of remapping to transfer velocity fields from unstructured \MPAS to \ROMS. 

It is important to bear in mind that neither finer \ROMS meshes nor higher-order interpolation schemes can remove the errors imposed by a coarse \MPAS representation in regions that require higher fidelity for better boundary forcing, definition of land masks, or estuarine structures. We have performed the resolution verification study to identify optimal characteristic length resolution, which can guide future efforts to improve resolution in regions of interest when designing CVT refinements for the \MPAS global ocean model.

\TODOc{add more relevant discussions from real world one-way runs; can reference videos of tracer evolution and discuss the chesapeake case and how that was achieved. Complications in getting that to work and why this augments global models.}

\TODOc{Placeholder for graph neural network exploration for volumetric remapping with conservation and monotonicity constratints}
%
